\documentclass[11pt]{article}
\usepackage[margin=1in]{geometry}
\usepackage{amsmath,amssymb,amsthm}
\usepackage{booktabs}
\usepackage{hyperref}

\newtheorem{theorem}{Theorem}
\newtheorem{lemma}[theorem]{Lemma}
\newtheorem{fact}[theorem]{Fact}
\theoremstyle{definition}
\newtheorem{convention}[theorem]{Convention}

\title{Disproof of TLMC Conjecture 00000000591\\
\large v2: the limit statement, formalized over an explicit infinite family}
\author{}
\date{October 2, 2026}

\begin{document}
\maketitle

\begin{abstract}
The new-direction conjecture attached to Wilf's conjecture predicts that on
numerical semigroups of embedding dimension $3$ the Wilf surplus ratio
$\bigl(e\cdot n(S)-g(S)-e\bigr)/g(S)$ tends to $0$ as $g(S)\to\infty$. We
disprove this with an explicit infinite family of complete-intersection
semigroups $S_k=\langle 6,\,2(6k{+}5),\,3(6k{+}5)\rangle$, $k=0,1,2,\dots$, for
which the ratio equals $\tfrac12-\tfrac{3}{2g_k}\ge \tfrac13$ for every $k$ and
converges to $\tfrac12$. The limit statement and its negation are formalized in
core Lean~4 (v4.33.1, no Mathlib) with zero axioms and zero \texttt{sorry}; the
main theorem is audited as \texttt{does not depend on any axioms}.
\end{abstract}

\section{The conjecture, quoted verbatim}

From \texttt{conjectures/00000000591.md}:

\begin{quote}
``Definition: Wilf's conjecture. Conjecture: $e\cdot n(S) \ge g(S) + e$ ($n(S)$
the number of nongaps, $g$ the Frobenius number). New-direction conjecture: on
embedding-dimension-3 semigroups the ratio of the surplus $e\cdot n(S) - g(S) - e$
to $g$ tends to $0$, with an explicit convergence rate.''
\end{quote}

\begin{convention}[Correspondence declaration]
$S$ ranges over numerical semigroups; $e=e(S)$ is the embedding dimension;
$g=g(S)$ is the \emph{Frobenius number} (the conjecture's notation; $g$ is not
a genus); $n=n(S)=\#\{s\in S : 0\le s\le g(S)\}$ is the number of nongaps in
$[0,g]$; $\mathrm{surplus}=e\cdot n(S)-g(S)-e$. Every semigroup constructed
below is of exactly this kind and is certified by exhaustive search, so no
object is substituted for the conjecture's own.
\end{convention}

\section{The explicit infinite family}

For $k=0,1,2,\dots$ put $r_k=6k+5$, which is coprime to $6$, and
\[
S_k=\langle 6,\;2r_k,\;3r_k\rangle=\langle pq,\;pr,\;qr\rangle,
\qquad (p,q,r)=(2,\,3,\,6k+5).
\]
Since $p,q,r$ are pairwise coprime, no generator lies in the subsemigroup
generated by the other two, so $S_k$ has embedding dimension exactly $3$; $S_k$
is a complete intersection (iterated gluing of $\langle 2\rangle$, $\langle 3\rangle$,
$\langle r_k\rangle$), hence symmetric. For this shape the Frobenius number is
\[
g_k \;=\; 2pqr-pq-pr-qr \;=\; 42k+29,
\]
and symmetry gives exactly $n_k=(g_k+1)/2=21k+15$ nongaps in $[0,g_k]$. The
Wilf surplus at $e=3$ is therefore
\[
\mathrm{surplus}_k = 3n_k-g_k-3 = 21k+13,
\qquad
\frac{\mathrm{surplus}_k}{g_k} = \frac{g_k-3}{2g_k} = \frac12-\frac{3}{2g_k}.
\]

\begin{theorem}[Disproof]\label{thm:main}
For every $k$, $g_k\ge 29>9$ and
$\mathrm{surplus}_k/g_k\ge 1/3$; moreover $\mathrm{surplus}_k/g_k\to 1/2$.
Hence there are embedding-dimension-3 numerical semigroups with arbitrarily
large $g$ and surplus ratio bounded below by the positive constant $1/3$: the
ratio does not tend to $0$. Conjecture $00000000591$ (new-direction part) is
false.
\end{theorem}

\begin{proof}
$g_k = 42k+29 \ge 29$ and
$\mathrm{surplus}_k/g_k-1/3 = \bigl(3(21k+13)-(42k+29)\bigr)/(3\,g_k)
= (21k+10)/(3\,g_k) \ge 0$. The limit is
$\lim_k (42k+26)/(84k+58) = 1/2$.
\end{proof}

\begin{fact}
The fraction $(\mathrm{surplus}_k, g_k)$ is in lowest terms for every $k$: a
common divisor divides $2\cdot\mathrm{surplus}_k-g_k=3$, and
$\mathrm{surplus}_k=21k+13\equiv 1 \pmod 3$.
\end{fact}

\section{The Lean formalization}

The Lean project (\texttt{lean4/}, core Lean v4.33.1, no Mathlib) proves the
discretized negation of the vanishing claim. With
$\texttt{third}=1/3\in\mathbb{Q}$, $\texttt{qzero}=0\in\mathbb{Q}$,
$\texttt{surplus}\ k = 21k+13$ and $\texttt{ratioDen}\ k = 42k+29$:

\begin{verbatim}
theorem main :
    exist epsilon : Rat, qzero < epsilon /\ forall N : Nat, exist k : Nat, N <= k /\
      epsilon.den * ratioNum k >= epsilon.num.natAbs * ratioDen k
\end{verbatim}

$\varepsilon$'s numerator is positive, so $\varepsilon.\mathrm{num.natAbs} =
\varepsilon.\mathrm{num} = 1$; both denominators are positive
($\texttt{third.den}=3$, $g_k\ge 29$), so
$\varepsilon.\mathrm{den}\cdot\mathrm{surplus}_k \ge
\varepsilon.\mathrm{num}\cdot g_k$ is exactly the cross-multiplied form of
$\varepsilon \le \mathrm{surplus}_k/g_k$: for every $N$ the member $k=N$
violates any decay of the ratio towards $0$. The witness is
$\varepsilon = 1/3$.

Supporting formalized facts, all audited as \texttt{does not depend on any
axioms} (\texttt{Check.lean}, 32 theorems; no \texttt{sorry}, no
\texttt{native\_decide}):

\begin{itemize}
\item \texttt{exact\_identity}: $2\cdot\mathrm{surplus}_k\cdot g_k = (g_k-3)\cdot g_k$,
i.e.\ $\mathrm{surplus}_k/g_k = (g_k-3)/(2g_k)$; \texttt{frob\_ge\_nine}: $g_k\ge 9$;
\texttt{wilf\_relation}: $3n_k = g_k + \mathrm{surplus}_k + 3$.
\item \texttt{zero\_lt\_third}: $0 < 1/3$ in $\mathbb{Q}$;
\texttt{third\_le\_ratio\_zero}: $1/3 \le 13/29$ and
\texttt{third\_le\_ratio\_one}: $1/3 \le 34/71$ \emph{as rational-number order
statements}, by definitional reduction of the rational order on ground values.
\item Anchoring certificates that the closed formulas describe the actual
semigroups (no object substitution): for $S_0=\langle 6,10,15\rangle$ an
exhaustive search certifies $29$ is a gap while $30,\dots,35$ are nongaps
(hence $29$ is exactly the Frobenius number) and that the nongap count on
$[0,29]$ is exactly $15$; likewise $S_1=\langle 6,22,33\rangle$ with $g=71$,
$n=36$.
\end{itemize}

On the shape of the symbolic inequality: core Lean's order on $\mathbb{Q}$
reduces to boolean cross-multiplication of the \texttt{num}/\texttt{den} fields,
and core \texttt{Rat} literals carry \texttt{propext}/\texttt{Quot.sound}
through their \texttt{by decide} coprimality fields; \texttt{omega} and
\texttt{simp} are avoided for the same reason. The main theorem therefore uses a
ground-constructed $\varepsilon=1/3$ and rational $0$, and states the symbolic
member inequality in the equivalent exact cross-multiplied $\mathbb{N}$ form;
the $\mathbb{Q}$-order form itself is proved on ground members.

\section{Numerical verification}

\texttt{reproduce.py} verifies by exhaustive enumeration (no closed formula is
trusted) that for every $k \le 300$ and spot checks up to $k=100{,}000$:
$g_k=42k+29$ is exactly the Frobenius number of $S_k$ (gap at $g_k$, the next
$6$ integers are nongaps), $n_k=21k+15$ is the exact nongap count on
$[0,g_k]$, $\mathrm{surplus}_k=21k+13$, $\gcd(\mathrm{surplus}_k,g_k)=1$, and
$\mathrm{surplus}_k/g_k\ge 1/3$. The prime-triple family
$\langle pq,pr,qr\rangle$ satisfies the same closed formulas:

\begin{center}
\begin{tabular}{lrrrrr}
\toprule
$(p,q,r)$ & $S_k$ & $g$ & $n$ & surplus & surplus$/g$\\
\midrule
$(2,3,5)$ & $\langle 6,10,15\rangle$ & 29 & 15 & 13 & $13/29 \approx 0.4483$\\
$(2,3,7)$ & $\langle 6,14,21\rangle$ & 43 & 22 & 20 & $20/43 \approx 0.4651$\\
$(2,5,7)$ & $\langle 10,14,35\rangle$ & 81 & 41 & 39 & $39/81 \approx 0.4815$\\
$(3,5,7)$ & $\langle 15,21,35\rangle$ & 139 & 70 & 68 & $68/139 \approx 0.4892$\\
\bottomrule
\end{tabular}
\end{center}

All ratios are $\ge 1/3$ and the family values converge to $1/2$
(e.g.\ $k=10^6$: $21000013/42000029 = 0.49999996$).

\section{What remains true}

Wilf's inequality $e\cdot n(S)\ge g(S)+e$ holds on every member here (e.g.\
$3\cdot 15 = 45 \ge 29+3$); only the attached vanishing-rate direction is
false. Moreover the pairing bound $n\le (g+1)/2$ gives
$\mathrm{surplus}/g \le (g-3)/(2g) < 1/2$ for every numerical semigroup of
embedding dimension $3$, so the symmetric family attains the exact asymptotic
ceiling: the conjectured limit $0$ fails in the strongest possible way.

\end{document}
